\chapter{3D and Video}

So far, we have explored how neural networks process static 2D images. However, the real world is three-dimensional and dynamic: objects have depth, volume, and move over time. Capturing and understanding this spatial and temporal richness is essential for advanced tasks such as 3D reconstruction, scene perception, object tracking, and action analysis in video.

In this chapter, we will discuss how deep learning extends beyond flat images to work with 3D data and temporal sequences. We will look at common representations for three-dimensional data (point clouds, meshes, voxels) and how network architectures adapt to these formats. We will also explore how to process video with networks that model both spatial and temporal dimensions.

We will review current challenges when working with 3D data and video sequences, such as the scarcity of labeled data and high computational requirements, as well as emerging trends to address them.

This chapter serves as a bridge between 2D perception and a richer understanding of space and time, laying the foundation for computer vision systems capable of interacting with complex and dynamic environments.



\section{3D Representations}

\subsection{Depth Map}

A \textbf{depth map} is a 2D representation in which each pixel stores information about the distance between the camera and objects in the scene in the real world. Instead of encoding color or light intensity, as in an RGB image, a depth map encodes depth: \textbf{small values represent points close to the camera, and large values indicate points farther away}. Visually, depth maps are usually represented in grayscale: white colors mean closer points and darker colors mean more distant points.

This representation is key to many computer vision and 3D graphics tasks, such as scene reconstruction, object pose estimation, autonomous navigation, and augmented reality. Figure \ref{depth_map} shows an example of two images and its depth maps. Note that in this case the depth map is no in grayscale, but a heatmap. In any case, there is only one value (one channel) for each pixel.

\begin{figure}[H]
\centering
\includegraphics[scale=0.8]{images/depth_map}
\caption{Example of two images and its depth maps.}
\label{depth_map}
\end{figure}

\subsection{Surface Normals}

A \textbf{surface normal} is a vector that is perpendicular to a surface at a given point. In computer vision and 3D graphics, a \textbf{normal map} is a 2D representation where each pixel stores the \textbf{orientation of the surface} at that point, encoded as a 3D vector $(x, y, z)$.

Unlike a depth map, which indicates how far a point is from the camera, a normal map describes how the surface is oriented in space. This is essential for capturing details of the local geometry of objects or simulating how light interacts with surfaces (shading and rendering).

Visually, normal maps are usually represented by assigning the $(x, y, z)$ components to RGB channels: red for the $X$ axis, green for the $Y$ axis and blue for the $Z$ axis. This allows surface orientations to be visualized with colors, although these do not have a real chromatic meaning but rather a geometric one.

\begin{figure}[H]
\centering
\includegraphics[scale=0.8]{images/surface_normals}
\caption{Example of two images and its normal maps.}
\label{surface_normals}
\end{figure}

\subsection{Implicit Functions}

\begin{figure}[H]
\centering
\includegraphics[scale=0.25]{images/imp_f}
\caption{Implicit function for a rabbit. Note that in the surface the function is zero, outside greater than zero, and inside smaller than zero.}
\end{figure}

An \textbf{implicit function} in 3D is a mathematical function $f(x, y, z)$ that defines a surface or volume continuously in space. The surface is defined by the set of points where the function vanishes, that is: $f(x, y, z)=0$. For example, a sphere of radius $r$ centered at $(0, 0, 0)$ can be implicitly defined as $f(x, y, z)=x^2+y^2+z^2-r^2=0$. Unlike explicit representations such as meshes, an implicit function does not directly store the geometry, but rather describes its shape using a continuous function. We can interpret $f(x, y, z)$ as the \textbf{distance to the surface} of the volume. Therefore, in a point outside the volume we have that $f(x, y, z) > 0$; in a point inside the volume we have that $f(x, y, z) < 0$; and in a point exactly in the surface of the volume we have that $f(x, y, z) = 0$.

Some advantages are the continuous and infinite representation (not limited to a discrete grid, you can evaluate any point); and the capability of capturing fine details with few parameters. However, some challenges still present are that it is expensive to evaluate the function to reconstruct dense geometry if we want to evaluate all points of a zone, so it requires specialized sampling and optimization techniques.

\subsection{Voxels}

A \textbf{voxel} (volumetric pixel or volume element) is the basic unit in a regular three-dimensional representation, analogous to a pixel in 2D. While a pixel represents a point in a flat image, a voxel represents a \textbf{small cube} in 3D space.

\begin{figure}[H]
\centering
\includegraphics[scale=0.8]{images/voxels}
\caption{Example of a chair and its voxel.}
\end{figure}

A voxel-based structure divides a three-dimensional volume into a regular grid of cubes, and each voxel can store information such as occupancy (is there matter in that cube?); color or intensity; or physical properties (density, temperature, etc.). Voxels are useful for representing 3D geometry. Usually they do not store information, only if there is a cube or not in that position, so its shape is $V\times V\times V$, but if they store some information, for example density and temperature at each point, the shape would be $V\times V\times V \times 3$ (there is a pixel or not, density and temperature).

They are a regular and direct representation of space. However, they are spatial inefficient: occupying an entire regular grid is very memory-intensive, especially at high resolution.

\subsection{Point Cloud}

A \textbf{point cloud} is a representation of a 3D shape or scene as a \textbf{set of points} in space. Each point has $(x, y, z)$ coordinates that define its position in a three-dimensional reference system, and sometimes includes additional attributes such as color (RGB), intensity, or surface normals.

Unlike a regular representation such as voxels, a point cloud is a \textbf{sparse} representation: it only stores information at the locations where data actually exists, making it efficient for describing 3D geometry.

They are lightweight and efficient because they only store relevant data. They are also fexible as the do not impose a fixed grid like voxels. However, they are unstructured (points have no explicit connectivity) and it is difficult to use in standard convolutional networks, as there is no regular grid topology.

\begin{figure}[H]
\centering
\includegraphics[scale=0.65]{images/cp}
\caption{Two different views for the reconstructed point cloud of an airplane.}
\end{figure}

\section{3D Convolutions}

\subsection{Model}

A \textbf{3D convolution} is an extension of the standard 2D convolution used in images, designed to process data that has three spatial or spatiotemporal dimensions. While in 2D, convolution covers the channels ($C$), the width ($W$) and height ($H$) of an image, in 3D it also considers \textbf{depth} ($D$), so the input shape is $(D, C, H, W$), where $D$ can be the spatial depth or the temporal axis in a sequence of images or video (frames during time). Figure \ref{3d_conv} gives an intuition of how these convolutions work.

3D convolutions learn features that combine spatial and temporal information and they are ideal for data with a regular 3D structure, but they have high computational and memory costs: kernels and activations grow with the third dimension.

The \textbf{kernel} of each channel is a 3D block with shape $(K_D, K_H, K_W)$ that slides in three dimensions and learns to capture volumetric or spatio-temporal patterns. Let $I$ be the input of the convolution and $O$ its output, then for a given channel:
$$O_{d, h, w}=\sum_{i=1}^{K_d}\sum_{j=1}^{K_h}\sum_{k=1}^{K_w}I_{d+i, h+j, w+k} \cdot K_{i, j, k}.$$
Therefore, the output shape is $(D_\text{out}, H_\text{out}, W_\text{out})$, where
$$D_\text{out}=\left\lfloor \dfrac{D+2P_D-D_D(K_D-1)-1}{S_D}+1\right\rfloor,$$
$$H_\text{out}=\left\lfloor \dfrac{H+2P_H-D_H(K_H-1)-1}{S_H}+1\right\rfloor,$$
$$W_\text{out}=\left\lfloor \dfrac{W+2P_W-D_W(K_W-1)-1}{S_W}+1\right\rfloor,$$
and the stride is $S=(S_D, S_H, S_W)$, the padding $P=(P_D, P_H, P_W)$ and the dilation $(D_D, D_H, D_W)$.

\textbf{Stride}, \textbf{padding} and \textbf{dilation} work the same as in 2D convolutions, but now they also include the depth dimension.

\begin{figure}[H]
\centering
\includegraphics[scale=0.65]{images/3d_conv}
\caption{Example of a 3D convolution for a specific channel.}
\label{3d_conv}
\end{figure}

\subsection{Max Pooling (3D) and Average Pooling (3D)}

\textbf{Max Pooling 3D} takes a block of shape $(P_D, P_H, P_W)$ as input and returns the maximum value of that block. \textbf{Average Pooling 3D} works the same, but returns the mean value of that block.

The output shape of these operations are the same as in the 3D convolution. Note that the max pooling is like if we had a kernel with zeros in all positions, except where the maximum value lives, where there would be a one; and in the average pooling is like if we had a $1/(K_DK_HK_W)$ in all positions.

\begin{figure}[H]
\centering
\includegraphics[scale=0.7]{images/maxp}
\caption{Max Pooling 3D. Note that the average pooling would be the same, but taking the mean of the block instead of the maximum.}
\end{figure}

\section{Predicting Representations}

\subsection{Predicting Depth Maps}

One straight way to train a model to predict depth maps is to train a fully convolutional network, as in the previous chapter, but now using only one output channel, as we only have to predict one value (depth) per pixel. We can use any metric that works for regression for the loss function, for example $L_2$ distance (mean squared error). This process is shown in Figure \ref{train_depth}.

\begin{figure}[H]
\centering
\includegraphics[scale=0.3]{images/train_depth}
\caption{Schema of the training for predicting depth maps.}
\label{train_depth}
\end{figure}

However, small, close objects look exactly he same as larger, farther away objects, so absolute scale or depth is ambiguous for a single image. You can see this with the example of Figure \ref{scale}. Therefore, we can improve the previous approach.

The \textbf{scale-invariant loss} is a loss function designed to make the model insensitive to absolute scales in the data. It is widely used, for example, in tasks such as depth estimation or 3D reconstruction, where the important thing is the relative structure (the relationships between points) and not the absolute magnitude of the distances. A typical way of defining it is:
$$\mathcal{L}(y, \hat{y})=\dfrac{1}{HW}\sum_{i=1}^H\sum_{j=1}^W\left(\log(y_{ij})-\log(\hat{y}_{ij})\right)^2 -  \dfrac{\lambda}{(HW)^2}\left(\sum_{i=1}^H\sum_{j=1}^W\left(\log(y_{ij})-\log(\hat{y}_{ij})\right)\right)^2,$$ where $y_{ij}$ and $\hat{y}_{ij}$ are the ground truth and prediction for pixel $(i, j$) and $0\leq\lambda\leq 1$ is a parameter typically set to one. The intuition behind this loss is that the first part measures the mean squared logarithmic error and the second part compensates for constant shifts in scale. If $\lambda = 0$, then we do not ``care'' about scale; if $0<\lambda<1$, we obtain a compromise; and if $\lambda = 1$ we obtain a total scale-invariant loss. This means that when $\lambda = 1$, if you multiply all predictions by a constant factor, i.e., $\hat{y}'_{ij}=c\hat{y}_{ij} ~ \forall i=1,\ldots,H; ~ \forall j=1,\ldots, W; ~ c>0$, the loss does not change: $\mathcal{L}(y, \hat{y}) = \mathcal{L}(y, c\hat{y})$. Note that if we do not use the logarithm in both terms, this property would not hold.

\begin{figure}[H]
\centering
\includegraphics[scale=0.5]{images/scale}
\caption{Problem: scale/depth ambiguity.}
\label{scale}
\end{figure}

\subsection{Predicting Surface Normals}

The approach to predicting surface normals is very similar to the one predicting depth maps. In fact, the process is all the same except that the output channels of the fully convolutional network are three (for the normal vector in each pixel) and the loss function used. Now we use as loss function the \textbf{cosine distance}:
$$\mathcal{L}(y, \hat{y})=\dfrac{1}{HW}\sum_{i=1}^H\sum_{j=1}^W\underbrace{\left(1-\dfrac{\vec{y}_{ij}\cdot \vec{\hat{y}}_{ij}}{||\vec{y}_{ij}||_2 ||\vec{\hat{y}}_{ij}||_2}\right)}_{\text{cosine distance}},$$
where $\vec{y}_{ij}=\left(y_{ij}^1, y_{ij}^2, y_{ij}^3\right)$ and $\vec{\hat{y}}_{ij}=\left(\hat{y}_{ij}^1, \hat{y}_{ij}^2, \hat{y}_{ij}^3\right)$ are the ground truth and predicted surface normals. Recall that we want the ground truth and predicted vectors to be similar, and the smaller the cosine distance, the closer these vectors are, as shown in Figure \ref{cosine}. More precisely this distance is in the range $[0, 2]$. If the vectors point in the same direction, the cosine distance is 0; if the vectors are orthogonal, the cosine distance is 1; and if the vectors are perfectly opposite, the cosine distance is 2. The rest of the cases are between these values.

\begin{figure}[H]
\centering
\includegraphics[scale=0.6]{images/cosine}
\caption{Cosine distance between two vectors. In this case, it will be near to zero.}
\label{cosine}
\end{figure}

\subsection{Simulating Implicit Functions}

To train a model that can simulate an implicit function, we can sample some random points of the space and use as labels their distance to the surface. Then, we can train any model with the inputs (points sampled) and labels generated. Note that the inputs will have dimension 3 and the output dimension 1.

\subsection{Predicting Voxels Classes}

If we want to predict to what class a voxel belongs to, we can apply a 3D convolutional neural network, flatten the resulting block and then use a multilayer perceptron (MLP) to obtain the class. The loss function will be cross entropy for multiclass and binary cross entropy for binary classification. An example of an architecture for this problem is shown in Figure \ref{voxel_class}.

\begin{figure}[H]
\centering
\includegraphics[scale=0.6]{images/voxel_class}
\caption{Architecture to predict what class a voxel belongs to. Note that we could also add 3D max pooling or average pooling between the 3D convolutions.}
\label{voxel_class}
\end{figure}

\subsection{Predicting Point Clouds}

To predict to what class a point cloud belongs to, we can pass each point through a MLP, aggregate all these output vectors with a \textbf{permutation invariant function}\footnote{A permutation invariant function is a function that for every possible permutation of the input, it returns the same value. The mean would be an example: $f(1, 2, 3)=(1+2+3)/3=6=f(1, 3, 2)=f(2, 1, 3)=f(2, 3, 1)=f(3, 1, 2)=f(3, 2, 1)$.}, for example mean or max along the $Y$ axis, and finally pass the obtained vector through another MLP to obtain the logits for each class. The loss function will be cross entropy for multiclass and binary cross entropy for binary classification. An example of an architecture for this problem is shown in Figure \ref{pc_class}.

\begin{figure}[H]
\centering
\includegraphics[scale=0.55]{images/pc_class}
\caption{Example of architecture to predict what class a point cloud belongs to.}
\label{pc_class}
\end{figure}

\section{Predicting Shapes}

\subsection{Metrics}

\subsubsection{Intersection over Union (IoU)}

This metric is the same seen for object detection. Just to recall, the IoU for two sets is
$$d_\text{IoU}(P, Q)=\dfrac{|P\cap Q|}{|P\cup Q|}.$$
However, in 3D shapes this metric is not the best one. It can not capture thin structures, it can only be applied to voxels and, in general, low values are not very meaningful.

The explanation about why it can only be applied to voxels is simple: voxels can only take a finite number of values. However, in other representations, like point clouds, the points can take continuous values. Therefore, the intersection between two sets of points will be almost always zero.

\subsubsection{Chamfer Distance}

The Chamfer distance is a widely used metric for comparing two point clouds ($P$ and $Q$) in 3D because it does not depend on exact correspondences between points. The mathematical formula is:
$$d_\text{C}(P, Q)=\dfrac{1}{|P|}\sum_{p\in P}\min_{q\in Q}\{||p-q||_2^2\}+\dfrac{1}{|Q|}\sum_{q\in Q}\min_{p\in P}\{||q-p||_2^2\}.$$

The intuition here is:
\begin{enumerate}
\item For each point $p\in P$, find the closest point in $Q$ and measure the squared euclidean distance.
\item Then do the same for each point $q\in Q$, finding its closest point in $P$.
\item Add the two means.
\end{enumerate}
In other words: it measures how well $P$ covers $Q$ and vice versa. The only problem is that in some cases it can be sensitive to outliers.

In addition, it has the following properties:
\begin{itemize}
\item It does not need to know which point in $P$ corresponds to which in $Q$.
\item The Chamfer distance is 0 for equal sets, i.e., if $P=Q$ then $d_\text{C}(P, Q)=0$.
\item It is symmetric: $d_\text{C}(P, Q)=d_\text{C}(Q, P)$.
\item It is differentiable, so you can use it as a loss function to train neural networks.
\item It can be applied to any shape representation.
\end{itemize}

\subsubsection{F1 Score}

In the context of 3D shapes (point clouds, meshes, voxels, etc.), F1 score is often used to evaluate the comparison between two 3D shapes. Let's say the set of predicted points is $P$ and the set of ground truth points is $G$. Then:
\begin{enumerate}
\item A distance threshold $\delta$ is set, for example $1\%$ of the object size.
\item A point $p\in P$ is true positive (TP) if it is within a distance less or equal than $\delta$ of a $G$ point. Otherwise, it is a false positive (FP).
\item A point $g\in G$ that is not covered by any $P$ point is considered as false negative (FN).
\item We define precision as the percentage of predicted points that are close to the actual shape:
$$\text{precision} = \dfrac{\text{TP}}{\text{TP}+\text{FP}}.$$
\item We define recall as the percentage of actual shape points that were covered by the prediction:
$$\text{recall} = \dfrac{\text{TP}}{\text{TP}+\text{FN}}.$$
\item Finally, we define:
$$\text{F1} = 2\cdot \dfrac{\text{precision}\cdot \text{recall}}{\text{precision}+ \text{recall}}.$$
\end{enumerate}
It can be applied to any shape representation, but it is sensitive to the threshold chosen. Nevertheless, it is believed it is the best shape prediction metric in common use.

\begin{figure}[H]
\centering
\includegraphics[scale=0.55]{images/f1}
\caption{F1 score for a predicted set of points and its grount truth set of points (``@t'' refers to the threshold used).}
\label{f1_example}
\end{figure}

\subsection{3D Datasets}

\textbf{ShapeNet} is one of the most widely used databases in the field of computer vision and 3D modeling. It contains more than 50k 3D models hierarchically organized into 55 categories, such as chairs, airplanes, and cars. These models are represented as polygonal meshes and have been carefully curated and aligned to facilitate tasks such as shape reconstruction, classification, and segmentation. ShapeNet also includes variants such as ShapeNetCore, which offers a cleaner and more consistent subset of objects, and ShapeNetSem, which provides detailed semantic annotations for each model. Thanks to its diversity and size, it has become a standard for training and evaluating neural networks in tasks related to understanding three-dimensional geometry.

\begin{figure}[H]
\centering
\includegraphics[scale=0.6]{images/datasets}
\caption{Some samples of the ShapeNet and Pix3D datasets.}
\end{figure}

\textbf{Pix3D}, on the other hand, is specifically designed to study the relationship between 2D images and 3D models. This dataset offers more than 10k real-world images of objects with their corresponding precisely aligned 3D models and associated pose (camera) data. Each instance includes mask and keypoint annotations, enabling tasks such as pose estimation, 3D shape retrieval from images, and novel view generation. Unlike ShapeNet, which focuses on synthetic data, Pix3D incorporates real-world complexity, such as occlusion, variable lighting, and visual noise, making it an essential resource for validating algorithms that must operate under non-ideal conditions.

\subsection{Generating Point Clouds}

The objective in this section is to obtain the point cloud of a 2D image. For that, we first use a CNN to extract features and reduce the size of the original image. Then we use two parallel branches:
\begin{itemize}
\item The first branch uses a fully connected network to predict $P_1$ points (each of these points has 3 elements $(x, y, z)$). The idea here is to capture the global shape of the object.
\item The second branch uses a CNN to obtain local points for each spatial position of the feature map extracted, concretely $P_2$ points per pixel. The idea here is to capture fine details and local variations. Note that the output of the CNN has to be reshaped from $3P_2 \times H'' \times W''$ to $H''W''P_2\times 3$.
\end{itemize}
The points from both branches are concatenated to obtain a tensor with shape $(P_1+H''W''P_2)\times 3$, that is, $P_1+H''W''P_2$ points. This results in the complete point cloud, which combines the global structure and local details. Finally, the Chamfer distance of the ground truth set of points and the predicted set of points is used as loss function.

\begin{figure}[H]
\centering
\includegraphics[scale=0.37]{images/pc_gen}
\caption{Generating a cloud point of a table. Red points represent the global shape and blue points represent local variations.}
\end{figure}


\section{Video}

\subsection{Introduction to Video}

In the context of computer vision, video analysis poses a number of additional challenges compared to processing still images. While an image can be represented as a two-dimensional array of pixels (height $\times$ width, with a possible third channel for color), a video introduces a temporal dimension, resulting in a three-dimensional (time $\times$ height $\times$ width) or even four-dimensional data volume if color channels are considered (time $\times$ channels $\times$ height $\times$ width). You can see this visually in Figure \ref{frames}. This temporal dimension is critical because it contains information about motion, scene dynamics, and the evolution of objects over time.

\begin{figure}[H]
\centering
\includegraphics[scale=0.095]{images/frames}
\caption{Videos incorporate a new dimension to images: the number of frames or time, usually denoted as $T$.}
\label{frames}
\end{figure}

Working with video involves addressing specific problems such as lighting variability throughout sequences, motion blur, partial or total occlusion of objects, and the need to capture temporal dependencies at different scales. Computationally, videos are much heavier than images: each second of video at 30 fps requires processing 30 images, which drastically increases storage, memory, and computing requirements.

From a deep learning perspective, this type of data requires models capable of capturing not only spatial information (patterns within each frame) but also temporal relationships between consecutive frames.

\subsection{Action Recognition}

Action recognition is one of the fundamental tasks in computer vision applied to video. It involves identifying and classifying actions or activities that occur in a temporal sequence, such as ``walking'', ``running'', ``jumping'', or ``clapping'' (Figure \ref{act_rec}). Unlike object recognition in static images, this task requires understanding not only the spatial appearance of individual frames but also the temporal evolution of movements.

One of the main difficulties in action recognition lies in the high intra-class variability: the same action can be executed in very different ways depending on the person, speed, camera angle, or lighting conditions. Furthermore, there are classes of actions with very subtle differences (e.g., ``drinking'' vs. ``eating'') that require models capable of capturing fine details in space and time.

Action recognition has direct applications in fields such as intelligent surveillance, human-computer interaction, augmented reality, and sports analysis, making it an area of great interest both academically and industrially. In terms of deep learning, various approaches have been proposed to address this task, which we will see along this section.

\begin{figure}[H]
\centering
\includegraphics[scale=0.5]{images/act_rec}
\caption{Difference between image classification and video classification.}
\label{act_rec}
\end{figure}

In many cases, as the videos are very heavy in memory (for example, for 30 fps HD $1920\times 1090$ is $\sim 10$ GB per minute), the following strategy is followed: the video is resized to a smaller shape and also the fps are reduced.

\subsubsection{Single-Frame CNN}

A simple idea is to train a normal CNN to classify images. Then we will use this CNN to classify each frame (or a subset of all the frames) of the video. Finally, we can take the mean to obtain the final prediction. The layer were the predictions are aggregated is usually called fusion layer. You can see the worflow in Figure \ref{single_frame}.

Although this technique may seem basic, it is a very good starting point when classifying videos and we can use it as a baseline to compare other methods.

\begin{figure}[H]
\centering
\includegraphics[scale=0.15]{images/single_frame}
\caption{Workflow of the single-frame CNN.}
\label{single_frame}
\end{figure}

\subsubsection{Late Fusion}

A more refined way to approach video classification consists on running a CNN for each frame to obtain their feature maps, then flatten these feature maps, and use a MLP to obtain the final class scores. This is called late fusion with FC layers.

\begin{figure}[H]
\centering
\includegraphics[scale=0.5]{images/late_fusion}
\caption{Schema of late fusion with FC layers.}
\end{figure}

Another variant of this method is late fusion with pooling, which also runs a CNN in each frame, obtains their feature maps, and instead of flatten and concatenate them, the mean is taken over the time (first dimension) and space (third and fourth dimensions). Finally, we use again a MLP to obtain the final class scores.

The main idea in late fusion is to get a high-level appearance of each frame and combine them. However, it is hard to compare low-level motion between frames. Note that it is called late fusion because the frames are merged after using the CNN.

\subsubsection{Early Fusion}

Early fusion tries to combine all frame in one tensor and use this as input for the CNN. That is, if the video has shape $T\times 3\times H\times W$, we reshape it into a tensor of shape $3T \times H\times W$ and then we fed this into a 2D CNN. The architecture of the CNN here is very flexible: we can use as many layers as we want and include a MLP at the end to predict the logits for each class.

One problem here is that the input dimension may be too large. Note that is called early fusion because the frames are merged before using the CNN.

\begin{figure}[H]
\centering
\includegraphics[scale=0.5]{images/early_fusion}
\caption{Schema of early fusion.}
\end{figure}

\subsubsection{3D CNN}

3D CNNs (seen previously in this chapter) apply convolutions to 4D tensors to fuse temporal information over the course of the network. Some 3D convolutional layers or 3D max/average pooling can be stacked. Finally, the feature maps are flattened and a MLP is used to classify the videeo.

Note that the process is very similar to early fusion, but using a 3D CNN instead of a 2D CNN, that it, the tensors have dimension $T\times C\times H \times W$ instead of $CT\times H \times W$.

\textbf{Temporal shift invariance} is a desirable property in models that process temporal sequences (such as video) and means that the model produces the same output even if the input is slightly shifted in time (for example, if the action occurs a little earlier or later).
\begin{itemize}
\item In the context of early fusion, data from different frames is combined before processing and it is not invariant to temporal shifts because the model sees the frames in a fixed order. If you change the order or shift the frames, the output changes.
\item 3D CNNs apply convolutions in space and time. If well designed (for example, with 3D kernels and without excessive temporal padding), they can capture local temporal patterns and be more robust to small shifts in time. However, they are not perfectly invariant, only tolerant of certain small shifts.
\end{itemize}

One of the most famous architectures is \textbf{C3D}, known also as the VGG of 3D CNNs. Its architecture is shown in Figure \ref{c3d}. The main problem is that it is very expensive computationally. Just to have some numbers to compare, AlexNet has $0.7$ GFLOP, VGG-16 $13.6$ and C3D $39.5$ GFLOP ($\sim 3\times$ VGG)!

\begin{figure}[H]
\centering
\includegraphics[scale=0.8]{images/c3d}
\caption{Architecture of C3D.}
\label{c3d}
\end{figure}

\textbf{Observation}: The single-frame CNN can handle videos of arbitrary length, while the videos for early fusion, late fusion and 3D CNNs must always have the same number of frames.

\subsubsection{Two-Stream Flow Networks}

\textbf{Optical flow} is a technique that calculates the apparent motion of pixels between two consecutive frames in a video. It represents motion vectors: how much and in what direction each pixel moves. The result is typically an image where each pixel has a direction and magnitude of motion. It captures temporal information (motion), which a normal 2D network (which sees static images) cannot directly learn. An example is shown in Figure \ref{optical_flow}.

\begin{figure}[H]
\centering
\includegraphics[scale=0.12]{images/optical_flow}
\caption{Example of optical flow. Note that $F_t(x, y)=(dx_t, dy_t)$ and $I_{t+1}(x, y)=I_t(x, y)+F_t(x, y)=(x_t+dx_t, y_t+dy_t)$.}
\label{optical_flow}
\end{figure}

There are some techniques to calculate optical flow, but they will not be discussed here. Some examples are TV-L1 or RAFT, and some libraries have them already implemented, like \texttt{cv2}.

\textbf{Two-Stream network} is an architecture proposed for action recognition in video. It has two branches (streams):
\begin{itemize}
\item Spatial stream (RGB): Takes normal RGB images (appearance) and learns what is seen (people, objects).
\item Temporal stream (optical flow): Takes optical flow maps (motion between frames) and learns how things move.
\end{itemize}

Both networks are CNNs (usually 2D), and are combined at the end to make the final prediction. They explicitly combine appearance and motion, usually outperform RGB-only CNNs in action tasks and require less data than 3D models like C3D, because it is based on pre-trained architectures (such as VGG or ResNet). However, optical flow is expensive to compute (offline or with a separate model), it is not easily trained end-to-end (the flow is not differentiable), double memory is used: two complete networks (more weight and computation) and they may be sensitive to errors in optical flow.

There are many ways to pass the inputs. Note that if we have $T$ frames, as for computing the optical flow between two images we need that two images, we can obtain at most $T-1$ optical flows (1 and , 2 and 3, \ldots, $T-1$ and $T$).
\begin{itemize}
\item For the spatial stream, we can use any architecture seen before (single-frame CNN, late fusion, early fusion or 3D CNN). Another approach to reduce the computation is to just pass one representative frame of the video.
\item For the temporal stream, the optical flows obtained are concatenated in the channels dimension (therefore, we obtain channels $dx_1, dy_1, \ldots, dx_{T-1}, dy_{T-1}$) and the input dimensions for the CNN are $2(T-1) \times H \times W$.
\end{itemize}
At the end, we can concatenate the feature maps obtained in both networks and use a MLP to obtain the final prediction or predict in each network and then average the results. Other option, in order to reduce the memory used, is to just use one CNN for the spatial and temporal streams; however, in this case, we can not use 3D CNN. A possible worflow of the model is shown in Figure \ref{two_stream}. Note that there are more possible options with this type of models. Not only those we have discussed in these notes, but also more in the literature.

\begin{figure}[H]
    \centering
    \includegraphics[scale=0.25]{images/excalidraw/two_stream.excalidraw.png}
    \caption{Possible workflow for two-stream flow networks.}
    \label{two_stream}
\end{figure}

\subsubsection{Inflated CNNs}

Inflated CNNs (I3D) adapt 2D convolutional networks (used in images) to the video (spatio-temporal) domain, inflating 2D convolutions to 3D convolutions. ``Inflate'' means extending 2D filters (height $\times$ width) to 3D (time $\times$ height $\times$ width). For example, a 2D kernel of size $3\times 3 \times 3$ is inflated to a 3D kernel of size $T\times 3 \times 3	\times 3$. Weights are repeated along the temporal dimension to capture motion between frames.

Let's suppose we have a 2D kernel $W_{\text{2D}}\in \mathbb{R}^{H \times W}$ and we want to inflate it to a 3D kernel $W_{\text{3D}}\in \mathbb{R}^{T\times H \times W}$. Then we can define:
$$W_{\text{3D}}[t, :, :] = \dfrac{1}{T}W_{\text{2D}} ~ \forall t\in\{1, \ldots, T\}.$$
The factor $1/T$ is optional, but it ensures that the magnitude of the inflated filter remains similar to the original (so as not to break the activation when going from 2D to 3D). This helps to use 2D weigths pretrained and helps the initial training. We can use this to do the training for a 3D CNN from a (usually) better starting point.

\subsection{Comparation of the Models}

The \textbf{Sports-1M} dataset, developed by Google Research in 2014, is one of the pioneering and most influential datasets for video analysis in the context of sports activity recognition. It consists of approximately 1.1 million videos automatically extracted from YouTube, which were classified into approximately 487 different sports-related categories. Unlike other manually annotated datasets, the labels in Sports-1M were automatically generated from the metadata, titles, and descriptions associated with the videos, which facilitates large-scale collection but introduces a considerable level of noise into the labels.

The videos vary widely in length, resolution, and visual quality, reflecting the diversity of online content. On average, each video is about five minutes long and can be labeled with multiple classes, making it a multi-class, multi-label classification dataset. This variability and volume make Sports-1M an ideal resource for testing deep learning models, especially those seeking to capture complex temporal patterns, such as 3D convolutional networks or LSTM-like recurrent architectures.

One of the highlights of Sports-1M is its relevance for video action and activity classification tasks, although its practical use presents certain challenges. Access to the dataset is not direct: due to licensing restrictions imposed by YouTube, the dataset does not include the videos themselves, but only a list of IDs that the user must manually download, assuming the videos are still available on the platform. Furthermore, automatic labeling can lead to incorrect or ambiguous classifications, so many researchers choose to perform additional filtering or validation before using the data in supervised tasks.

Some examples of the dataset are shown in Figure \ref{sports}. Also, the models seen in the previous section have been tested in this dataset. You can see these comparations in Figure \ref{comp1}, where the metric used is top-5 accuracy\footnote{Top-5 accuracy is a metric primarily used in multiclass classification, especially in tasks with a large number of classes, to assess whether the model has included the correct class among its five most likely predictions. Rather than requiring the model to accurately predict the most likely class (as in top-1 accuracy), top-5 accuracy allows a prediction to be considered correct if the true class is among the five most likely classes assigned by the model.}.

\begin{figure}[H]
\centering
\includegraphics[scale=0.21]{images/sports}
\caption{Some examples of the Sports-1M dataset.}
\label{sports}
\end{figure}

\begin{figure}[H]
\centering
\includegraphics[scale=0.4]{images/comp1}
\caption{Top-5 accuracy on Sports-1M for different models.}
\label{comp1}
\end{figure}



\subsection{Long-Term Temporal Structure}

So far, all our temporal CNNs only model local motion between frames in very short clips of around 3 seconds. What about long-term structure?

\subsubsection{LSTMs for Video}

A natural way of modelling sequences is using recurrent networks, specially Long Short-Term Memory networks (LSTMs). To apply LSTMs to video, you can treat each video as a sequence of frames (images). First, you extract features from each frame using a convolutional network (CNN), such as ResNet or VGG. Then, those features—vectors that summarize the visual content of each image—are passed sequentially to an LSTM, which learns the temporal dependencies between frames (e.g., motion or an action that develops over time). Finally, the output of the LSTM is used in a fully connected layer to obtain the final prediction. Note that the main idea here is to extract local information with the CNNs and then obtain a global view of the video with the LSTM.

The CNN used can be trainable, but usually its weights are fixed to make training faster. Also, the output can be one prediction per frame or just one final prediction, as in Figure \ref{lstm_video}. This depends on the task we are doing. For example, for video classification only one output is needed, but for video segmentation or if we know more than one action will be done and we want to classify all actions done in the video we probably want one output per frame.



\begin{figure}[H]
\centering
\includegraphics[scale=0.3]{images/lstm_video}
\caption{Using a LSTM for video.}
\label{lstm_video}
\end{figure}




\subsubsection{Recurrent Convolutional Neural Networks}

A \textbf{Recurrent Convolutional Neural Network (RCNN)} combines convolutional networks (CNNs) with recurrent networks (RNNs, such as LSTMs or GRUs) to leverage both spatial and temporal or sequential information. The CNN is first used to extract spatial features from inputs, such as images or video frames. The CNN captures local patterns such as edges, shapes, or textures. Then, these features are passed to the RNNs (or LSTMs), which model the sequence—that is, obtain how those features change over time or in a defined order.LSTMs for video are a particular case of this architecture.

A \textbf{Multi-Layer RCNN} is simply a RCNN with several recurrent layers stacked on top of each other, allowing it to learn deeper representations in both space and time.
\begin{itemize}
\item First recurrent-convolutional layer: Processes inputs (such as frame feature maps) by capturing basic patterns in space and time.
\item Additional layers: Each higher layer receives the output of the previous one and learns more abstract or higher-level representations.
\item Final output: After several RCNN layers, a dense or classification layer can be added to produce the final output (e.g., an action label in a video).
\end{itemize}

\begin{figure}[H]
\centering
\includegraphics[scale=0.5]{images/lstm_2}
\caption{General view of the multi-layer RCNN.}
\label{lstm_2}
\end{figure}

\begin{figure}[H]
\centering
\includegraphics[scale=0.35]{images/lstm_3}
\caption{The features for layer $L$ and timestep $t$ are computed based on the features from layer $L$ and timestep $t-1$ and the features from layer $L-1$ and timestep $t$.}
\label{lstm_3}
\end{figure}

Having multiple layers provide greater modeling capabilities: they capture more complex relationships in the sequence; pattern hierarchy: lower layers detect local or simple movements and higher layers understand more complex dynamics (e.g., an entire action); and more temporal context: by stacking layers, you expand the range of temporal information the model can consider. Figures \ref{lstm_2} and \ref{lstm_3} provide a general view of the model and how the features in each block are computed, respectively.



\subsubsection{SlowFast Networks}

\textbf{SlowFast networks} are a neural network architecture for action recognition in video, proposed by Facebook AI Research. Its main idea is to process video in two parallel temporal streams:
\begin{itemize}
\item \textbf{Slow pathway}: Processes fewer frames over time (low sampling rate); focuses on rich semantic information (such as stable objects or scenes); and uses deep 3D convolutions, but with a low temporal frequency.
\item \textbf{Fast pathway}: Processes many more frames per second, but with a lighter network; captures rapid dynamics, such as subtle movements or gestures; and has lower spatial representation capacity, but higher temporal resolution.
\end{itemize}
Both paths are periodically merged (e.g., through concatenation or summation) to combine slow context with fast detail, improving the model's ability to understand complex actions that combine motion and context.
There are also two key parameters:
\begin{itemize}
\item Time speedup factor ($\alpha$): The fast path processes $\alpha$ times more frames than the slow path. A typical value is $\alpha=8$.
\item Channel reduction factor ($\beta$): The fast path has $1/\beta$ times fewer channels than the slow path. This makes it computationally lighter. Two typical values are $\beta=1/6$ and $\beta=1/8$.
\end{itemize}

The fast pathway is composed of a relatively small channel capacity compared to slow network. This design method increases computational efficiency with a small channel capacity. With this structure, it is relatively difficult to obtain spatial information of a frame. However, this defect is compensated by the slow pathway. Indeed, the fast network focuses on the action capture with a small computation cost.

\begin{figure}[H]
\centering
\includegraphics[scale=0.29]{images/slowfast}
\caption{Architecture of a SlowFast network.}
\label{slowfast}
\end{figure}



\subsubsection{Comparation of the Models}

One comparation for some of the models we have seen done in the literature is using the dataset Kinetics-400. Kinetics-400 is a video human action recognition dataset developed by DeepMind. It contains 400 action classes, with at least 400 distinct clips per class, each lasting approximately 10 seconds, extracted from different YouTube videos In total, it includes more than 300k short videos. The results obtained are shown in Figure \ref{comp2}.

\begin{figure}[H]
\centering
\includegraphics[scale=0.6]{images/comp2}
\caption{Top-1 accuracy comparation for different models on Kinetics-400.}
\label{comp2}
\end{figure}


\subsection{Other Problems in Video}

During the chapter we have mainly focused on video classification because it is the most common task. However, there are some other problems that are worth mentioning:
\begin{itemize}
\item \textbf{Temporal Action Localization (TAL)}: The goal is to detect when an action occurs in a long, unsegmented video. The input is a continuous video (e.g., several minutes) and the output are time intervals [start, end] along with the action label. For example, in a video of a person training, detect that a ``jump rope'' occurs between seconds 15 and 22, and a ``push-up'' occurs between seconds 30 and 40. Some challenges are that actions can be brief and subtle and that there may be multiple actions in succession or overlapping.
\item \textbf{Spatio-Temporal Action Detection}: The goal is to detect when and where an action occurs in the video, that is, locate it in time and space (image). The input is a video and the output a bounding box around the actor and the action label for each frame. For example, to detect that a person is ``kicking a ball'' between seconds 10 and 13, and in each frame in that range, draw a box around the person. The main challenges are the challenges of TAL plus the challenges of spatial detection.
\end{itemize}

\begin{figure}[H]
\centering
\includegraphics[scale=0.45]{images/temporal_action_localization}
\caption{Example of temporal action localization.}
\label{temporal_action_localization}
\end{figure}

\begin{figure}[H]
\centering
\includegraphics[scale=0.5]{images/std}
\caption{Examples of spatio-temporal action detection.}
\label{std}
\end{figure}

\section{Exercises}

\setcounter{question}{0}

\begin{question}
Do by hand a 3D convolution over a video with just one channel. Let's say the input is $3\times 1\times 5\times 5$. Use a $2\times 2\times 2$ kernel with the values you want (I would suggest small integers). Use also stride 1, 0 padding and no dilation.
\end{question}

\begin{question}
Repeat the exercise before with stride 2 and padding 1.
\end{question}

\begin{question}
Write the code to apply a 3D convolution over a video with just one channel. You can assume stride is 1 and padding is 0. Check that the outputs are the same as in the first exercise.
\end{question}

\begin{question}
Extend the previous exercise with general padding and stride.
\end{question}

\begin{question}
Write the code for a 3D convolutional layer. The input parameters are kernel size, stride, padding, input channels and output channels. Now the input can have more than one channel. Check that the results are the same as in the \href{https://docs.pytorch.org/docs/stable/generated/torch.nn.Conv3d.html}{3D convolution of PyTorch}.
\end{question}

\begin{question}
Show that if you multiply all predictions by a constant factor, i.e., $\hat{y}'_{ij}=c\hat{y}_{ij} ~ \forall i=1,\ldots,H; ~ \forall j=1,\ldots, W; ~ c>0$, the scale invariant loss does not change, that is, $\mathcal{L}(y, \hat{y}) = \mathcal{L}(y, \hat{y}')$.
\end{question}

\begin{question}
Write the code of a model to predict depth maps. Look for a dataset in Internet. Use first the mean squared error and then the scale-invariant loss.
\end{question}

\begin{question}
Write the code of a model to predict surface normal. Look for a dataset in Internet.
\end{question}

\begin{question}
Write the code of a model to predict the class of a voxel. Look for a dataset in Internet.
\end{question}

\begin{question}
Write the code of a model to predict to which class a point cloud belongs to. Look for a dataset in Internet.
\end{question}

\begin{question}
Show that:
\begin{itemize}
\item The Chamfer distance is zero for two equal sets.
\item The Chamfer is symmetric.
\end{itemize}
\end{question}

\begin{question}
Write a program to calculate the F1 score given two set of points. Remember the $\delta$ parameter. Check that the code is correct using an example similar to Figure \ref{f1_example}.
\end{question}

\begin{question}
Write the code for video classification using different models:
\begin{itemize}
\item Single-Frame CNN.
\item Late Fusion.
\item 3D CNN.
\item Recurrent Convolutional Neural Network.
\item SlowFast Network.
\end{itemize}
Compare and comment the results obtained.
\end{question}

